% docs/papers/theory.tex
%
% Univalence Gravity: A Constructive Formalization of Discrete
% Entanglement-Geometry Duality in Cubical Agda
%
% Draft paper — machine-checked results only, no conjectures.
% Every theorem cited here type-checks in Agda 2.8.0 + agda/cubical.

\documentclass[11pt,a4paper]{article}

% ── Packages ────────────────────────────────────────────────────────
\usepackage[utf8]{inputenc}
\usepackage[T1]{fontenc}
\usepackage{lmodern}
\usepackage{amsmath,amssymb,amsthm}
\usepackage{mathtools}
\usepackage{hyperref}
\usepackage{cleveref}
\usepackage{booktabs}
\usepackage{enumitem}
\usepackage{xcolor}
\usepackage[margin=2.5cm]{geometry}
\usepackage{microtype}

% ── Theorem environments ────────────────────────────────────────────
\newtheorem{theorem}{Theorem}[section]
\newtheorem{lemma}[theorem]{Lemma}
\newtheorem{proposition}[theorem]{Proposition}
\newtheorem{corollary}[theorem]{Corollary}
\newtheorem{definition}[theorem]{Definition}
\newtheorem{remark}[theorem]{Remark}

% ── Custom commands ─────────────────────────────────────────────────
\newcommand{\N}{\mathbb{N}}
\newcommand{\Z}{\mathbb{Z}}
\newcommand{\Qi}{\mathbb{Z}[i]}
\newcommand{\Qge}{\mathbb{Q}_{\geq 0}}
\newcommand{\Qten}{\mathbb{Q}_{10}}
\newcommand{\refl}{\mathsf{refl}}
\newcommand{\ua}{\mathsf{ua}}
\newcommand{\transport}{\mathsf{transport}}
\newcommand{\funExt}{\mathsf{funExt}}
\newcommand{\isSet}{\mathsf{isSet}}
\newcommand{\isContr}{\mathsf{isContr}}
\newcommand{\isProp}{\mathsf{isProp}}
\newcommand{\Type}{\mathsf{Type}}
\newcommand{\Scut}{S_{\mathrm{cut}}}
\newcommand{\Lmin}{L_{\mathrm{min}}}
\newcommand{\Agda}{\textsc{Agda}}
\newcommand{\module}[1]{\texttt{#1}}

% ── Title ───────────────────────────────────────────────────────────

\title{%
  Univalence Gravity:\\[4pt]
  A Constructive Formalization of Discrete\\
  Entanglement-Geometry Duality in Cubical Agda%
}

\author{Sven Bichtemann}

\date{Draft — \today}

% ════════════════════════════════════════════════════════════════════
\begin{document}
\maketitle

% ════════════════════════════════════════════════════════════════════
\begin{abstract}
We present a machine-checked formalization in Cubical Agda of five
pillars of a discrete holographic universe---geometry, causality,
gauge matter, curvature, and quantum superposition---coexisting in
a single formal artifact verified by computational transport along
Univalence paths.

The central results, all constructively proven:
\begin{enumerate}[nosep]
\item A \emph{generic} discrete Ryu--Takayanagi theorem
  ($\Scut = \Lmin$ on every boundary region) parameterized by
  an abstract patch interface, instantiated on twelve patches
  spanning 1D trees, 2D pentagonal tilings, and 3D cubic
  honeycombs.
\item A discrete Gauss--Bonnet theorem
  ($\sum \kappa(v) = \chi(K) = 1$) for both hyperbolic and
  spherical pentagonal tilings, discharged by $\refl$.
\item A sharp discrete Bekenstein--Hawking bound
  $S(A) \leq \mathrm{area}(A)/2$ with $1/(4G) = 1/2$ in
  bond-dimension-1 units---an exact rational verified by
  $\refl$ on closed $\N$ terms, eliminating the
  constructive-reals obstacle for the entropy-area relationship.
\item Structural acyclicity (no closed timelike curves) as a
  consequence of $\N$ well-foundedness, enforced by the type
  signature rather than by a separate axiom.
\item A five-line amplitude-polymorphic quantum bridge theorem
  lifting the per-microstate correspondence to arbitrary finite
  superpositions.
\end{enumerate}

The proof architecture introduces three engineering innovations:
a \emph{generic bridge pattern} (one proof, $N$ instances),
\emph{orbit reduction} (717 regions $\to$ 8 proof cases), and
the \emph{abstract barrier} for propositional proofs sealed against
downstream RAM cascades.  No axioms are postulated; all transport
computes via the $\ua\beta$ reduction rule of Cubical Type Theory.
\end{abstract}

\tableofcontents

% ════════════════════════════════════════════════════════════════════
\section{Introduction}
\label{sec:intro}

The holographic principle~\cite{Maldacena1997,Ryu2006} asserts that
the physics of a bulk gravitational spacetime is encoded on its
lower-dimensional boundary.  The Ryu--Takayanagi (RT) formula
\[
  S(A) \;=\; \frac{\mathrm{Area}(\gamma_A)}{4 G_N}
\]
quantifies this duality: the entanglement entropy $S(A)$ of a
boundary region $A$ equals the area of the minimal bulk surface
$\gamma_A$ homologous to $A$, measured in Planck units.

This paper presents the first machine-checked formal artifact in
which the combinatorial core of this correspondence---together with
discrete curvature, causal structure, gauge matter, and quantum
superposition---is stated, proven, and computed within a single
constructive type theory.

\paragraph{Contributions.}
\begin{enumerate}[nosep]
\item A generic bridge theorem (\cref{sec:bridge}) producing a
  full enriched type equivalence + Univalence path + verified
  transport for \emph{any} finite patch satisfying a four-field
  abstract interface.
\item A sharp half-bound (\cref{sec:bekenstein}) identifying the
  discrete Newton's constant as $1/(4G) = 1/2$ by $\refl$.
\item A curvature-agnostic bridge (\cref{sec:wick}) that is
  literally the same Agda term for both AdS-like and dS-like
  geometries.
\item Three proof-engineering patterns (\cref{sec:engineering})
  that scale the formalization from 8-region toy models to
  1885-region exponentially growing hyperbolic patches.
\end{enumerate}

\paragraph{Scope and honesty.}
We prove properties of \emph{finite, discrete, combinatorial}
structures.  We do \emph{not} prove convergence to smooth
geometry, connection to continuous gauge groups, or that
``transport along a Univalence path'' has physical meaning.
Four hard walls remain between this formalization and continuum
physics (\cref{sec:walls}).

% ════════════════════════════════════════════════════════════════════
\section{Background: Cubical Agda and Univalence}
\label{sec:background}

We work in Cubical Agda~\cite{VezzosiMortbergAbel2021}, an
implementation of Cubical Type Theory~\cite{CCHM2018} in which:
\begin{itemize}[nosep]
\item Paths are functions from the abstract interval $I$ with
  endpoints $i_0, i_1$.
\item Function extensionality ($\funExt$) is a theorem, not an axiom.
\item The Univalence Axiom ($\ua : A \simeq B \to A \equiv B$) is
  implemented via Glue types and has a computation rule:
  \[
    \ua\beta(e, a) \;:\;
    \transport(\ua\, e)\, a \;\equiv\; \mathrm{equivFun}\, e\, a
  \]
  Transport along a $\ua$ path \emph{computes} as the forward map
  of the equivalence.
\end{itemize}

\begin{definition}[Scalar types]
\label{def:scalars}
The nonneg.\ scalar type is $\Qge = \N$ (\module{Util/Scalars}).
The signed curvature type is $\Qten = \Z$, where integer $n$
represents $n/10$ (\module{Util/Rationals}).  All arithmetic
reduces judgmentally on closed constructor terms, enabling every
agreement and Gauss--Bonnet proof to be discharged by $\refl$.
\end{definition}

% ════════════════════════════════════════════════════════════════════
\section{The Generic Holographic Bridge}
\label{sec:bridge}

\subsection{The PatchData Interface}

\begin{definition}[\module{Bridge/GenericBridge}]
\label{def:patchdata}
A \textbf{PatchData} record consists of:
\begin{align*}
  &\mathsf{RegionTy} : \Type_0 \\
  &S_\partial : \mathsf{RegionTy} \to \Qge \\
  &L_B : \mathsf{RegionTy} \to \Qge \\
  &\mathsf{obs\text{-}path} : S_\partial \equiv L_B
\end{align*}
Nothing about pentagons, cubes, curvature, dimension, gauge
groups, or Coxeter reflections appears.
\end{definition}

\subsection{The Generic Enriched Equivalence}

\begin{theorem}[Generic Bridge, \module{Bridge/GenericBridge}]
\label{thm:generic-bridge}
For any $\mathit{pd} : \mathsf{PatchData}$, the parameterized
module $\mathsf{GenericEnriched}(\mathit{pd})$ produces:
\begin{enumerate}[nosep]
\item Enriched types
  $\mathsf{EnrichedBdy} = \Sigma_{f} (f \equiv S_\partial)$,
  $\mathsf{EnrichedBulk} = \Sigma_{f} (f \equiv L_B)$
  (reversed singletons, contractible).
\item $\mathsf{enriched\text{-}equiv} :
  \mathsf{EnrichedBdy} \simeq \mathsf{EnrichedBulk}$.
\item $\mathsf{enriched\text{-}ua\text{-}path} :
  \mathsf{EnrichedBdy} \equiv \mathsf{EnrichedBulk}$.
\item $\mathsf{enriched\text{-}transport} :
  \transport(\mathsf{enriched\text{-}ua\text{-}path})\,
  \mathsf{bdy} \equiv \mathsf{bulk}$.
\item $\mathsf{abstract\text{-}bridge\text{-}witness} :
  \mathsf{BridgeWitness}$.
\end{enumerate}
The transport step closes by contractibility of
$\Sigma_{x \in A}(x \equiv a)$ via $\isContr \to \isProp$,
then $\ua\beta$ reduces transport to the forward map.
\end{theorem}

\subsection{Orbit Reduction}

\begin{definition}[\module{Bridge/GenericBridge}]
An \textbf{OrbitReducedPatch} carries a large $\mathsf{RegionTy}$,
a small $\mathsf{OrbitTy}$, a classification function
$\mathsf{classify} : \mathsf{RegionTy} \to \mathsf{OrbitTy}$, and
orbit-level observables $S_{\mathrm{rep}}, L_{\mathrm{rep}}$ with
pointwise $\refl$ agreement.  The PatchData is extracted
automatically:
\[
  S_\partial(r) = S_{\mathrm{rep}}(\mathsf{classify}\, r), \qquad
  \mathsf{obs\text{-}path} =
  \funExt(\lambda r.\, \mathsf{rep\text{-}agree}(\mathsf{classify}\, r))
\]
The one-function composition
$\mathsf{orbit\text{-}bridge\text{-}witness} :
\mathsf{OrbitReducedPatch} \to \mathsf{BridgeWitness}$
produces the full proof-carrying result.
\end{definition}

Twelve patch instances have been verified, from 8-region trees to
1885-region $\{5,4\}$ depth-7 patches (942$\times$ orbit reduction).

% ════════════════════════════════════════════════════════════════════
\section{Discrete Gauss--Bonnet}
\label{sec:gauss-bonnet}

\begin{theorem}[Discrete Gauss--Bonnet, \module{Bulk/GaussBonnet}]
\label{thm:gauss-bonnet}
For the 11-tile filled patch $K$ of the $\{5,4\}$ hyperbolic tiling
with combinatorial curvature
$\kappa(v) = 1 - \deg(v)/2 + \sum_{f \ni v} 1/\mathrm{sides}(f)$:
\[
  \sum_{v} \kappa(v) \;=\; \chi(K) \;=\; 1
\]
The proof is $\refl$: the class-weighted $\Z$ sum normalizes
judgmentally to $\mathsf{pos}\, 10 = \mathsf{one}_{10}$.
\end{theorem}

Five vertex classes suffice (Strategy~A from~\cite{HaPPYInstance}):
$5 \cdot (-2) + 5 \cdot (-1) + 5 \cdot (+2) + 5 \cdot (-1) +
10 \cdot (+2) = +10$ in tenths.

A parallel theorem holds for the $\{5,3\}$ spherical tiling with
\emph{positive} interior curvature $\kappa = +1/10$
(\module{Bulk/DeSitterGaussBonnet}).

% ════════════════════════════════════════════════════════════════════
\section{The Discrete Bekenstein--Hawking Bound}
\label{sec:bekenstein}

\begin{theorem}[Half-Bound, \module{Bridge/HalfBound}]
\label{thm:half-bound}
For any cell-aligned boundary region $A$ on a verified patch:
\[
  S(A) \;\leq\; \frac{\mathrm{area}(A)}{2}
\]
with a tight achiever where $2 \cdot S(A) = \mathrm{area}(A)$.
The discrete Newton's constant is $1/(4G) = 1/2$ in
bond-dimension-1 units.
\end{theorem}

\begin{proof}[Proof sketch]
Decompose $\mathrm{area}(A) = n_{\mathrm{cross}} + n_{\mathrm{bdy}}$
where $n_{\mathrm{cross}}$ counts bonds crossing from $A$ to its
complement and $n_{\mathrm{bdy}}$ counts boundary legs of $A$-cells.
The max-flow $S(A)$ is bounded by two independent cuts:
\begin{enumerate}[nosep]
\item Sever all crossing bonds: $S \leq n_{\mathrm{cross}}$.
\item Sever all $A$-boundary legs: $S \leq n_{\mathrm{bdy}}$.
\end{enumerate}
Therefore
$S \leq \min(n_{\mathrm{cross}}, n_{\mathrm{bdy}})
\leq (n_{\mathrm{cross}} + n_{\mathrm{bdy}})/2
= \mathrm{area}/2$.
The Agda formalization composes these via the
$\mathsf{two\text{-}le\text{-}sum}$ arithmetic lemma and the
$\mathsf{from\text{-}two\text{-}cuts}$ generic lemma.
Per-instance witnesses (717 for Dense-100, 1246 for Dense-200) are
generated by the Python oracle and sealed behind
$\mathsf{abstract}$.
\end{proof}

Numerically confirmed across 32,134 regions on 4 tilings
($\{4,3,5\}$, $\{5,4\}$, $\{4,4\}$, $\{5,3\}$), 4 growth
strategies, and 3 bond capacities---zero violations.

\begin{remark}[Elimination of the constructive-reals wall]
The original \emph{Entropic Convergence Conjecture} required a
$\mathsf{ConvergenceWitness}$ involving Cauchy completeness.
The sharp half-bound replaces this with
$\mathsf{HalfBoundWitness}$ requiring only $\N$ arithmetic and
$\refl$.  The discrete Newton's constant is an exact rational,
not a real-valued limit.
\end{remark}

% ════════════════════════════════════════════════════════════════════
\section{Discrete Wick Rotation}
\label{sec:wick}

\begin{theorem}[Curvature-Agnostic Bridge,
    \module{Bridge/WickRotation}]
\label{thm:wick}
The holographic bridge is curvature-agnostic: the same Agda term
serves both the negatively curved $\{5,4\}$ (AdS-like) and the
positively curved $\{5,3\}$ (dS-like) pentagonal tilings.  A
$\mathsf{WickRotationWitness}$ record packages:
\begin{enumerate}[nosep]
\item A shared $\mathsf{BridgeWitness}$ (produced by the generic
  theorem, never inspecting curvature).
\item An AdS Gauss--Bonnet witness ($\kappa_{\mathrm{int}} = -1/5$).
\item A dS Gauss--Bonnet witness ($\kappa_{\mathrm{int}} = +1/10$).
\item Euler coherence: $\chi_{10}^{\mathrm{AdS}} \equiv
  \chi_{10}^{\mathrm{dS}}$ by $\refl$.
\end{enumerate}
No complex numbers, no analytic continuation, no Lorentzian geometry.
\end{theorem}

% ════════════════════════════════════════════════════════════════════
\section{Causal Structure and No Closed Timelike Curves}
\label{sec:causal}

\begin{theorem}[NoCTC, \module{Causal/NoCTC}]
\label{thm:noctc}
In any time-stratified causal poset (where each causal link strictly
increases the $\N$-valued time index), no closed timelike curve
exists:
\[
  \mathsf{CausalChain}\; e\; e'
  \;\to\; \mathsf{CausalLink}\; e'\; e
  \;\to\; \bot
\]
\end{theorem}

\begin{proof}
$\mathsf{chain\text{-}step\text{-}<}$ composes single-step
orderings via $<_\N$-transitivity to get
$\mathsf{time}(e) <_\N \mathsf{time}(e)$; then
$<_\N$-irreflexivity (from $\mathsf{snotz}$ + $\mathsf{injSuc}$)
yields $\bot$.
\end{proof}

The acyclicity is a \emph{structural consequence of the type
signature} of $\mathsf{CausalLink}$---the type system prevents
time travel by construction.

% ════════════════════════════════════════════════════════════════════
\section{Gauge Theory: Matter as Topological Defects}
\label{sec:gauge}

\begin{definition}[\module{Gauge/Holonomy}]
A \textbf{ParticleDefect} at a face $f$ under gauge connection
$\omega$ is:
\[
  \mathsf{ParticleDefect}(\omega, \partial f)
  \;:=\;
  \neg\bigl(\mathsf{holonomy}(\omega, \partial f)
  \equiv \varepsilon\bigr)
\]
\end{definition}

\begin{theorem}[\module{Gauge/Holonomy}]
For the quaternion group $Q_8 \subset SU(2)$ on the 6-tile star
patch with connection $\mathsf{starQ8\text{-}i}$ (bond
$\mathsf{bCN0} = q_i$, rest $= q_1$), the holonomy at the central
face reduces to $q_i$ and $\mathsf{defect\text{-}Q8\text{-}i}$
inhabits $\mathsf{ParticleDefect}$.  Non-commutativity is witnessed:
swapping bond assignments changes the holonomy from $q_k$ to
$q_{-k}$.
\end{theorem}

The bridge is \textbf{gauge-agnostic}: the dimension functor
$\dim : \mathrm{Rep}\, G \to \N$ extracts scalar bond capacities,
and the generic theorem operates on these scalars without seeing
the group.

% ════════════════════════════════════════════════════════════════════
\section{Quantum Superposition Bridge}
\label{sec:quantum}

\begin{theorem}[Quantum Bridge, \module{Quantum/QuantumBridge}]
\label{thm:quantum}
For any finite superposition $\psi$ of configurations with
amplitudes in any algebra $A$, if $S(\omega) = L(\omega)$ for every
microstate $\omega$, then $\mathbb{E}[\psi, S] \equiv
\mathbb{E}[\psi, L]$.
\end{theorem}

\begin{proof}
Structural induction on the list $\psi$.  Base: $\refl$.
Step: $\mathsf{cong}_2\; {+_A}\;
(\mathsf{cong}\; (\alpha \cdot_A {-})\;
(\mathsf{cong}\; \mathsf{embed}_\N\; (\mathsf{eq}\; \omega)))\;
(\mathsf{IH})$.
Five lines.  No ring axioms.
\end{proof}

Verified with 3-configuration $Q_8$ superpositions exhibiting
destructive interference ($\Qi$ amplitudes), with the bridge
holding through the cancellation.

% ════════════════════════════════════════════════════════════════════
\section{Proof Engineering}
\label{sec:engineering}

\subsection{The Generic Bridge Pattern}
One parameterized proof ($\sim$30 lines) subsumes all twelve
bridge instances.  Adding a new patch requires zero new hand-written
Agda proof---only oracle-generated data modules.

\subsection{Orbit Reduction}
For the Dense-100 patch: 717 regions share 8 distinct min-cut
values.  The proof operates on the 8-constructor orbit type; a
1-line lifting connects to the full 717-constructor region type.
The depth-7 $\{5,4\}$ layer achieves 942$\times$ reduction (1885
regions, 2 orbits).

\subsection{The \texttt{abstract} Barrier}
Wrapping propositional proofs in $\mathsf{abstract}$ seals them
against downstream re-normalization, halting the RAM cascade
documented in Agda Issue \#4573.  Since $\_\leq_{\Qge}\_$ is
propositional ($\isProp$ by $\N$ cancellation + $\isSet_\N$),
sealing loses zero information.

\subsection{The Oracle Pipeline}
Seventeen Python scripts generate Agda modules with explicit
$\refl$-based witnesses.  The oracle is the search engine; Agda
is the checker.  This is the same division of labor as the Four
Color Theorem proof in Coq.

% ════════════════════════════════════════════════════════════════════
\section{The Schematic Tower and Convergence Certificate}
\label{sec:tower}

The $\mathsf{SchematicTower}$ assembles verified patches into a
resolution sequence with monotonicity witnesses:

\begin{center}
\begin{tabular}{lcccl}
\toprule
Level & Patch & Regions & Orbits & Max $S$ \\
\midrule
0 & Dense-50  & 139  & --- & 7 \\
1 & Dense-100 & 717  & 8   & 8 \\
2 & Dense-200 & 1246 & 9   & 9 \\
\bottomrule
\end{tabular}
\end{center}

Monotonicity: $7 \leq 8$ via $(1, \refl)$; $8 \leq 9$ via
$(1, \refl)$.  The $\mathsf{DiscreteBekensteinHawking}$ capstone
type carries half-bounds at each level:

\begin{align*}
  &\mathsf{DiscreteBekensteinHawking} : \Type_1 \\
  &\mathsf{DiscreteBekensteinHawking} =
    \mathsf{ConvergenceCertificate3L\text{-}HB}
\end{align*}

% ════════════════════════════════════════════════════════════════════
\section{The Hard Boundaries}
\label{sec:walls}

Four independent obstacles remain between the discrete formalization
and continuum physics:

\begin{enumerate}[nosep]
\item \textbf{Infinite-dimensional path integrals.}  The quantum
  bridge works for finite sums only.
\item \textbf{Continuous gauge groups.}  $Q_8 \subset SU(2)$ is
  the farthest reach; $|G| \to \infty$ limits require constructive
  Lie theory.
\item \textbf{Lorentzian signature.}  The causal poset provides
  time-directionality but no metric signature.
\item \textbf{Fermionic matter.}  The gauge construction produces
  bosonic defects; Grassmann variables are unexplored constructively.
\end{enumerate}

The constructive-reals wall for the entropy-area relationship is
\textbf{partially bypassed}: $1/(4G) = 1/2$ is an exact rational
verified by $\refl$, not a real-valued limit.

% ════════════════════════════════════════════════════════════════════
\section{Related Work}
\label{sec:related}

The HaPPY code~\cite{Pastawski2015} introduced the tensor-network
model of holographic quantum error correction on which this
formalization is based.  The Ryu--Takayanagi
formula~\cite{Ryu2006} and its covariant
generalization~\cite{Wall2012} motivate the min-cut/max-flow
observables.  Jacobson's thermodynamic derivation of Einstein's
equations~\cite{Jacobson1995} provides the physical context for
the entropy-area relationship.

In formal methods, the Cubical Agda
framework~\cite{VezzosiMortbergAbel2021} provides computational
Univalence; the oracle-plus-checker paradigm follows the Four Color
Theorem~\cite{Gonthier2005} and Kepler
Conjecture~\cite{Hales2017} proofs.

To our knowledge, no previous work has machine-checked a
holographic correspondence at the type-equivalence level in any
proof assistant, nor combined it with discrete curvature, causal
structure, gauge matter, and quantum superposition in a single
formal artifact.

% ════════════════════════════════════════════════════════════════════
\section{Conclusion}
\label{sec:conclusion}

We have constructed, from first principles, a discrete holographic
universe inside a proof assistant: boundary entanglement entropy
exactly equals bulk minimal surface area (RT), spacetime curvature
satisfies Gauss--Bonnet, the arrow of time is enforced by $\N$
well-foundedness, matter manifests as topological defects in $Q_8$
Wilson loops, quantum superposition preserves the correspondence,
and the sharp half-bound $S \leq \mathrm{area}/2$ identifies the
discrete Newton's constant as an exact rational.

The bridge is \emph{geometrically blind}: it depends on four inputs
(a region type, two observable functions, a path between them) and
never inspects dimension, curvature, gauge group, or tiling
geometry.  This blindness is the structural content of the claim
that the holographic correspondence is a property of information
flow, not of geometry.

Whatever the continuum theory turns out to be, it must be consistent
with the discrete theorems proven here.

% ════════════════════════════════════════════════════════════════════
\section*{Acknowledgements}

The author thanks Anthropic's Claude for serving as an invaluable
sounding board and pair-programmer throughout the mathematical
groundwork and architectural design phases.

% ════════════════════════════════════════════════════════════════════
\bibliographystyle{alpha}

\begin{thebibliography}{99}

\bibitem[CCHM18]{CCHM2018}
C.~Cohen, T.~Coquand, S.~Huber, A.~M\"ortberg.
\newblock Cubical Type Theory: A Constructive Interpretation of the
  Univalence Axiom.
\newblock \emph{Journal of Automated Reasoning}, 60(2):195--230, 2018.

\bibitem[Gon05]{Gonthier2005}
G.~Gonthier.
\newblock A computer-checked proof of the Four Colour Theorem.
\newblock Unpublished manuscript, Microsoft Research, 2005.

\bibitem[Hal+17]{Hales2017}
T.~Hales et~al.
\newblock A Formal Proof of the Kepler Conjecture.
\newblock \emph{Forum of Mathematics, Pi}, 5:e2, 2017.

\bibitem[Jac95]{Jacobson1995}
T.~Jacobson.
\newblock Thermodynamics of Spacetime: The Einstein Equation of State.
\newblock \emph{Physical Review Letters}, 75(7):1260--1263, 1995.

\bibitem[Mal97]{Maldacena1997}
J.~Maldacena.
\newblock The Large $N$ Limit of Superconformal Field Theories and
  Supergravity.
\newblock \emph{Advances in Theoretical and Mathematical Physics},
  2(2):231--252, 1998.

\bibitem[Pas+15]{Pastawski2015}
F.~Pastawski, B.~Yoshida, D.~Harlow, J.~Preskill.
\newblock Holographic Quantum Error-Correcting Codes: Toy Models for
  the Bulk/Boundary Correspondence.
\newblock \emph{Journal of High Energy Physics}, 2015(6):149, 2015.

\bibitem[RT06]{Ryu2006}
S.~Ryu, T.~Takayanagi.
\newblock Holographic Derivation of Entanglement Entropy from AdS/CFT.
\newblock \emph{Physical Review Letters}, 96(18):181602, 2006.

\bibitem[VMA21]{VezzosiMortbergAbel2021}
A.~Vezzosi, A.~M\"ortberg, A.~Abel.
\newblock Cubical Agda: A Dependently Typed Programming Language with
  Univalence and Higher Inductive Types.
\newblock \emph{Journal of Functional Programming}, 31:e8, 2021.

\bibitem[Wal12]{Wall2012}
A.~Wall.
\newblock Maximin Surfaces, and the Strong Subadditivity of the
  Covariant Holographic Entanglement Entropy.
\newblock \emph{Classical and Quantum Gravity}, 31(22):225007, 2014.

\bibitem[HaPPYInstance]{HaPPYInstance}
S.~Bichtemann.
\newblock HaPPY Instance specification.
\newblock \texttt{historical/development-docs/09-happy-instance.md}
  in the repository, 2025.

\end{thebibliography}

% ════════════════════════════════════════════════════════════════════
\appendix
\section{Module Index (Selected)}
\label{app:modules}

\begin{center}
\small
\begin{tabular}{ll}
\toprule
Module & Content \\
\midrule
\module{Bridge/GenericBridge} & PatchData, OrbitReducedPatch,
  GenericEnriched \\
\module{Bridge/HalfBound} & from-two-cuts, HalfBoundWitness \\
\module{Bridge/WickRotation} & WickRotationWitness \\
\module{Bridge/SchematicTower} & TowerLevel, DiscreteBekensteinHawking \\
\module{Bridge/BridgeWitness} & Universal result record \\
\module{Bulk/GaussBonnet} & $\sum\kappa = \chi$ by $\refl$ \\
\module{Causal/NoCTC} & Structural acyclicity \\
\module{Causal/LightCone} & FutureCone, Spacelike \\
\module{Gauge/Holonomy} & holonomy, ParticleDefect \\
\module{Gauge/Q8} & 400-case associativity by $\refl$ \\
\module{Quantum/QuantumBridge} & 5-line superposition bridge \\
\module{Quantum/AmplitudeAlg} & $\Qi$ instance \\
\bottomrule
\end{tabular}
\end{center}

\section{Verification Commands}
\label{app:verification}

\begin{verbatim}
agda src/Bridge/GenericBridge.agda     # Theorem 1 (RT)
agda src/Bulk/GaussBonnet.agda         # Theorem 2 (GB)
agda src/Bridge/HalfBound.agda         # Theorem 3 (BH)
agda src/Bridge/WickRotation.agda      # Theorem 4 (Wick)
agda src/Causal/NoCTC.agda             # Theorem 5 (NoCTC)
agda src/Gauge/Holonomy.agda           # Theorem 6 (Matter)
agda src/Quantum/QuantumBridge.agda    # Theorem 7 (Quantum)
\end{verbatim}

Each command loads the target module and all transitive dependencies.
The full repository type-checks via
\texttt{agda src/Bridge/SchematicTower.agda}.

\end{document}